\documentclass[10pt]{article}
\usepackage[margin=2.5cm]{geometry}
\usepackage{booktabs}
\usepackage{float}
\usepackage{graphicx}

\title{Labwork 4\\\large HPC -- 2D grid for grayscale}
\author{
Le Nghiem Thanh Thuy - 2540023\\
University of Science and Technology of Hanoi
}
\date{}

\begin{document}
\maketitle

\section{Preparation}
The image is the same as in labwork 3, \texttt{astronaut()} from \texttt{skimage.data},
size $512\times512$ with 3 channels. In labwork 3 I flattened it to $(262144, 3)$ and used
a 1D grid. Here the image keeps its shape $(512, 512, 3)$ and the threads are organised in
a 2D grid, so a thread finds its pixel with a row index and a column index instead of a
single index.

\section{Results}
I used the same method as in labwork 3: the image stays on the device so only the kernel
is measured, each configuration is launched once as a warm-up and then 100 times, and I
take the average.

\begin{table}[H]
\centering
\caption{Kernel time for the 2D block sizes (Tesla T4 on Google Colab).}
\label{tab:2d}
\begin{tabular}{lrrr}
\toprule
\textbf{Block size} & \textbf{Threads} & \textbf{Grid} & \textbf{Time (ms)} \\
\midrule
$(8, 8)$   &   64 & $(64, 64)$ & 0.0993 \\
$(16, 8)$  &  128 & $(32, 64)$ & 0.0469 \\
$(16, 16)$ &  256 & $(32, 32)$ & 0.0411 \\
$(32, 8)$  &  256 & $(16, 64)$ & 0.0402 \\
$(16, 32)$ &  512 & $(32, 16)$ & 0.0375 \\
$(32, 32)$ & 1024 & $(16, 16)$ & 0.0376 \\
$(64, 16)$ & 1024 & $(8, 32)$  & 0.0404 \\
\bottomrule
\end{tabular}
\end{table}

\begin{figure}[H]
\centering
\includegraphics[width=\textwidth]{blocksize2d.png}
\caption{Kernel time of the 1D version (blue) and the 2D version (orange).}
\label{fig:blocksize2d}
\end{figure}

The best 2D configuration is $(16, 32)$ with 0.0375 ms and the best 1D one is 256 threads and 512 threads
with 0.0607 ms.

\section{Comment}
The 2D grid is more convenient to write, do not have to reshape anything.


\end{document}
